  \def\stat{sorokin}



\def\tit{МЕТОД СЕГМЕНТАЦИИ ИЗОБРАЖЕНИЙ\\ НА~ОСНОВЕ 

КВАДРОДЕРЕВА}



\def\titkol{Метод сегментации изображений на основе квадродерева}





\def\aut{Ю.\,А.~Маньяков$^1$, А.\,И.~Сорокин$^2$}



\def\autkol{Ю.\,А.~Маньяков, А.\,И.~Сорокин}



\titel{\tit}{\aut}{\autkol}{\titkol}



\index{Маньяков Ю.\,А.}

\index{Сорокин А.\,И.}

\index{Maniakov Yu.\,A.}

\index{Sorokin A.\,I.}





%{\renewcommand{\thefootnote}{\fnsymbol{footnote}}

%\footnotetext[1]

%{Статья подготовлена при частичной поддержке РФФИ (проекты 18-29-03124-мк и~18-29-03081-мк).}} 





\renewcommand{\thefootnote}{\arabic{footnote}}

\footnotetext[1]{Орловский филиал Федерального исследовательского центра <<Информатика и управ\-ле\-ние>> Российской академии наук, \mbox{maniakov\_yuri@mail.ru}}

\footnotetext[2]{Орловский филиал Федерального исследовательского центра <<Информатика и управ\-ле\-ние>> Российской академии наук, \mbox{webdi@mail.ru}}





 \label{st\stat}



 \thispagestyle{headings}

 

 \vspace*{-12pt}  

  

  \Abst{Ппредставлен метод сегментации цветного изображения на основе квадродерева, 

состоящий из нескольких шагов. Первый из них~--- обнаружение границ, основанное на 

трехканальном цвете и двух масках. Далее применяется алгоритм утончения контура. Алгоритм 

сегментации делится на две части. В~первой части изображение разделяется на максимальное 

число сегментов. Во второй части сегменты объединяются на основе FSM-таб\-ли\-цы. 

В~статье представлены результаты сегментации цветных изображений, полученные на основе 

описанного алгоритма.}

  

  \KW{цветное изображение; сегментация; квадродерево; ребро; граница; истончение; 

уменьшение цвета; разделение; слияние; пиксель}

  

\DOI{10.14357/08696527200410}



\vskip 10pt plus 9pt minus 6pt



 \thispagestyle{headings}

 

 \vspace*{-12pt}

  

\section{Введение}



  Сегментация~---  одна из фундаментальных операций, используемых в 

обработке изображений и,~в~частности, в области компьютерного зрения~[1] 

для извлечения необходимой информации из сцены. Несмотря на то что этому 

вопросу посвящено немало исследований, задача остается актуальной~[2]. 

В~общем, можно сказать, что сегментация заключается в разбиении 

изображения на различные непересекающиеся области. Такие области, как 

правило, однородны по некоторому признаку, вычисленному на основе 

изображения.

  

  В основе алгоритмов сегментации лежат различные подходы~[3], из которых 

можно выделить четыре основных.

  \begin{enumerate}[1.]

  \item Обнаружение контуров. Этот подход широко используется для 

сегментации че\-рно-бе\-лых изображений~[4]. В~таких изображениях поиск 

краев осуществляется методом градиентной аппроксимации на основе 

вычисления первой производной (операторы Канни, Собеля и др.~[5]). С~точки 

зрения визуальной оценки они выглядят хорошо, но не соответствуют 

требованиям непрерывности и минимальной толщины контурных линий. 

  

  Градиентный поиск цветных изображений осуществляется, в основном, двумя 

способами~[6]: для всех вариаций цветовых каналов используется одна мера; 

градиент канала вычисляется отдельно и объединяет эти вычисления по 

определенным критериям.

  \item Выращивание регионов. Этот процесс начинается с отбора исходного 

<<семени>>. Точки вокруг семени группируют по некоторому критерию 

однородности. Для сегментации цветных изображений используется большое 

число методов, основанных на обработке полутоновых изображений~[7]. 

Необходимо учитывать, что при использовании данного подхода имеет место 

зависимость полученных результатов от порядка обработки точек изображения. 

Также в этом случае необходимо проанализировать каждый пиксель и~соотнести 

его с определенной однородной областью.

  \item Методы разделения и слияния. Эта группа использует изображение в 

качестве начальной области для разделения~\cite{3-sor}. После создания 

однородных областей деление заканчивается. Кроме того, соседние регионы 

связаны на основе квадратичного дерева. Связанные области должны быть 

гладкими и непрерывными. Для этого используют марковское случайное поле, 

стабилизированные уравнения обратной диффузии и др.

  \item Кластеризация. В процессе кластеризации пиксели изображения делятся 

на группы, которые не пересекаются. Характеристики объектов, входящих 

в~группу, должны быть близкими. В~противном случае объекты, принадлежащие 

к разным группам, должны существенно отличаться. Соответственно, 

расположение пикселей в каждом кластере будет определяться цветовыми 

вариациями изображения. В~этом случае задача кластеризации определяется как 

поиск определенного числа паттернов, однозначно определяющих число 

областей~[7].

  \end{enumerate}

  

  В данной работе предложен метод сегментации изображений на основе 

разделения и слияния. Этот метод состоит из поиска ребер; истончения 

вычисленных границ; уменьшения цвета; разбиения изображения квадрата на 

основе цвета и учета информации о границах; объединения однородных 

областей на основе графа. Данный алгоритм применяется для извлечения 

информации о сегментах и анализа их для поиска опорных точек в задачах 

стереозрения.

     

\section{Сегментация}



  В формальном виде сегментация изображения может быть представлена 

следующим образом~\cite{3-sor}. Пусть изображение~$I$ размером $W\times H$ 

является некоторой дискретной двумерной функцией $f(x,y)$, где~$I$~--- это двумерная мат\-ри\-ца, 

в~которой определена точка на изображении, а соответствующие 

точечные элементы матрицы задают ее цветовую характеристику.

  

  Пусть $P$~--- предикат однородности. Тогда сегментация разбивает~$I$ на 

множество~$N$~областей~$S_n$, $n\hm = 1,\ldots, N$, таких что

  \begin{enumerate}[(1)]

  \item $\bigcup\nolimits_{n=1}^N S_n=I$, $S_n\cap S_m=0$, $n\not=m$;

  \item $P(S_n)=1$ $\forall n$;

  \item $P(S_n\cup S_m)=0$, $\forall S_n, S_m$--- смежных.

  \end{enumerate}

  

  Первое и второе условия говорят о том, что каждый пиксель изображения 

принадлежит только одной области~\cite{8-sor}. В~третьем указывается, что 

области однородны.

   

  Процесс сегментации изображения состоит из нескольких этапов (рис.~1).

  

  \begin{figure} %fig1

\vspace*{1pt}

 \begin{center}

 \mbox{%

 \epsfxsize=116.949mm 

\epsfbox{man-1.eps}

 }

 \end{center}

\vspace*{-11pt}

  \Caption{Этапы сегментации}

  \end{figure}

  

  Первоначально осуществляется выделение контуров на основе цветовых 

значений. После этого выполняется утончение контура. Далее уменьшается 

число цветов на изображении и на основе информации о цвете и о контурах 

выполняется расщепление, а затем слияние.

  

  \textbf{Структура данных.} Как упоминалось ранее, сегментация основана на 

квадродереве~\cite{9-sor}. Изображение хранится в дереве таким образом, что 

каждый узел имеет четыре потомка~--- NW, NE, SW и SE.

  

  Таким образом, дескриптор прямоугольной области изображения и узла 

дерева ($\mathrm{Node}$) описывается следующим образом:

  $$

  \mathrm{Node}= \left\{ \mathrm{Level}, \left\{ \begin{matrix}

  \mathrm{RegonID}\\ \mathrm{PartitionID}\\ \mathrm{ChNum} \end{matrix}\right\},

  \mathrm{avgColor, nPoint, nBorder, Rect, Env}\right\}.

  $$

Здесь $\mathrm{Level}$~---  высота текущего узла квадродерева; $\mathrm{RegonID}$~--- уникальный 

идентификатор текущего узла, который вычисляется следующим 

образом~\cite{10-sor}:

\begin{equation}

\mathrm{RegonID= PartitionId \cdot 10 +ChNum}\,;

\label{e2-sor}

\end{equation}

$\mathrm{PartitionID}$~--- идентификатор родителя; $\mathrm{ChNum}$~--- номер узла в~текущем 

разделе: NW~--- 1, NE~--- 2, SW~--- 3, SE~--- 4; 

$\mathrm{avgColor}$~--- среднее значение 

цвета;

  $\mathrm{nPoint}$~--- число точек в текущем узле;

  $\mathrm{nBorder}$~--- число точек, принадлежащих контуру в текущем узле;

  $\mathrm{Rect}$~--- границы узла: {Left}, {Right}, {Top}, {Bottom}; 

$\mathrm{Env}$~--- идентификаторы соседних узлов: $\mathrm{NorthID}$, $\mathrm{EastID}$, $\mathrm{SouthID}$, 

$\mathrm{WestID}$.

  

  \textbf{Обнаружение контуров}. Для обнаружения контуров используется 

подход, описанный  

в~работе~\cite{11-sor}~--- поиск изменения цвета с помощью масок (рис.~2). Для 

снижения вычислительной нагрузки используются две маски: горизонтальная 

(0$^\circ$) и вертикальная (90$^\circ$) (рис.~2,\,\textit{б}).

  

\begin{figure} %fig2

\begin{center}

 \mbox{%

 \epsfxsize=86.699mm 

\epsfbox{man-2.eps}

 }

 \end{center}

\vspace*{-11pt}

\Caption{Применение масок: (\textit{а})~расположение пикселей: (\textit{б})~маски 0$^\circ$ 

и~90$^\circ$}

\end{figure}



  Пороговое значение вычисляется следующим образом:

  $$

  T=\left( \mathrm{dp}\,(i,j)\vert \mathrm{Avg}\right)\,;

  $$

  

  

  \noindent

  \begin{multline*}

  \left( \mathrm{dp}(i,j)\vert 0^\circ\right)= {}\\

  {}=\sum\limits_{i=0}^{\max I} 

\sum\limits_{j=0}^{\max J}\left[ \left\vert \left( 

 K_{\max} R(i,j-1) +G(i,j-1) +B(i,j-1)\right) %*\cdot 

 a(i,j-1)+{}\right.\right.\\

 \left.\left.{}+\left( K_{\max} R(i,j+1) +G(i,j-1)

   +B(i,j+1)\right) %* 

   a(i,j+1)\right\vert \right]\,;

  \end{multline*}

  

  \vspace*{-12pt}

  

  \noindent

  \begin{multline*}

    \left( \mathrm{dp}\,(i,j)\vert 90^\circ \right) ={}\\

    {}=\sum\limits_{i=0}^{\max I} 

\sum\limits_{j=0}^{\max J}\left[ \left\vert

  \left( K_{\max} R(i-1,j) +G(i-1,j) +B(i-1,j) \right) %*\cdot 

  b(i-1,j)+{}\right.\right.\\

\left.\left.  {}+\left( K_{\max} R(i+1,j)+ 

G(i+1,j)+B(i+1,j)\right) %* 

b(i+1,j)\right\vert \right]\,.

  \end{multline*}

  

  \textbf{Утончение контура}. Для прореживания границ применяется алгоритм, 

описанный в~\cite{12-sor}, в~котором применяются два шага итерации~--- один 

для горизонтальных и вертикальных линий, второй для наклонных. После этого 

необходимо провести отбор. Наиболее подходящие границы выбираются на 

основе компонентов связности. Учитываются непрерывность и длина границ. 

Результатом этой операции становится создание бинарной матрицы $\mathrm{brdMtx}$ 

размерности $W\times H$.

   

  \textbf{Прореживание цвета}. Для уменьшения цвета изображения 

используется подход, описанный  

в~\cite{13-sor}. Такой выбор обусловлен простотой реализации и~качеством 

получаемых результатов.

  

  \textbf{Расщепление}. Процесс расщепления осуществляется путем деления 

родительского узла на четыре дочерних. Для каждого квадранта проверяется, 

насколько однородна часть изображения, лежащая в его границах. В~общем 

случае, если предикат, примененный к четырем дочерним узлам, истинен, то 

расщепления не происходит:



\noindent

  $$

  \mathrm{Node}_n= \begin{cases}

  \bigcup\limits_{j=1}^4 \mathrm{Node}_{nj}, & \mbox{если } P(\mathrm{Node}_n)=\mathrm{false};\\

   \mathrm{Node}_n & \mbox{иначе}.

  \end{cases}

  $$

  

  Для каждого дочернего узла вычисляется среднее значение его цвета 

$\mathrm{avgColor}$ и число граничных пикселей $\mathrm{numBrd}$. Если манхэттенское 

расстояние между $\mathrm{avgColor}$ и цветом узла $\mathrm{pColor}$ меньше~$T_{\mathrm{clr}}$ или 

$\mathrm{numBrd} \hm\geq T_{\mathrm{brd}}$, то текущий узел делится на четыре.

  

  Процесс расщепления можно представить следующим образом:

\vspace*{4pt}



\noindent

  $$

  P\left( \mathrm{Node}_n\right) =

  \begin{cases}

  \mbox{false}, & \mbox{если } \displaystyle \sum\limits_{x=0}^H \sum\limits_{y=0}^W 

\left\vert \mathrm{avgColor}_{\mathrm{Node}_n}\left(f(x,y)\right) -{}\right.\\

&\hspace*{70pt}\left.{}-\mathrm{pColor}_{x,y\in \mathrm{Node}_n} 

\left(f(x,y)\right)\right\vert \leq T_{\mathrm{clr}}\,;\\[3pt]

  \mbox{false}, & \mbox{если } \mathrm{numBrd}_{\mathrm{Node}_n}\geq T_{\mathrm{brd}}\,;\\

  \mbox{true} & \mbox{иначе}.

  \end{cases}

  $$

\vspace*{-2pt}

      

  Алгоритм разделения узлов выглядит следующим образом. На вход подается 

изображение с уменьшенным числом цветов $\mathrm{Iq}$ и $\mathrm{brdMtx}$. В~цикле 

выполняются следующие действия:

  \begin{itemize}

\item вычисление координат потомка;

\item поиск среднего цвета для $\mathrm{Node}_n$ каждого канала:

\item  вычисление дистанции между $\mathrm{avgColor}_{\mathrm{Node}_n}$ и $\mathrm{pColor}_{x,y\in 

\mathrm{Node}_n}$;

\item если $\mathrm{dst}_c\leq T_{\mathrm{clr}}$ или $\mathrm{brd}_n\geq T_{\mathrm{brd}}$, то выполнение следующего 

действия, иначе конец;

\item добавление информации о текущем узле, вычисление $\mathrm{RegionID}$ по 

формуле~(\ref{e2-sor}).

\end{itemize}



  На выходе получаем $\bigcup\nolimits_{n=1}^N S_n\hm=I$.

  

  Результатом выполнения этого алгоритма будет максимальное разбиение 

изображений.

  

  \textbf{Слияние}. Этап объединения начинается с определения соседей для 

каждого узла. Эта операция основана на модификации подхода, предложенного 

в работе~\cite{14-sor}, т.\,е.\ используется конечный автомат (КА), 

представленный в таблице.



  В отличие от~\cite{14-sor} используются четыре состояния. Это делается для 

снижения вычислительной нагрузки.

  

  Два значения (число и тире) в ячейке указывают на то, что поиск соседей 

в~этом направлении должен продолжаться для следующего прохода (R, L, D и~U). 

Одно значение в ячейке означает, что текущий квадрант последний и~поиск 

в~этом направлении необходимо остановить.

  

%\begin{floatingfigure}{57mm}  

  \begin{table}

\small 

%\vspace*{-3pt}



  \begin{center}

  \begin{tabular}{|c|c|c|c|c|}

  \multicolumn{5}{c}{\tabcolsep=-3pt\begin{tabular}{l}

 Конечный автомат для соседей

 квад-\\ рантов по четырем направлениям\end{tabular}}\\

 \multicolumn{5}{c}{\ }\\[-6pt]

 \hline

\tabcolsep=-1pt\begin{tabular}{c}Направ-\end{tabular}&\multicolumn{4}{c|}{Квадрант}\\ 

\cline{2-5}

\multicolumn{1}{|c|}{ление} & 1& 2& 3& 4\\ 

\hline

&\multicolumn{1}{p{16pt}|}{\hspace*{16pt}}&

\multicolumn{1}{p{16pt}|}{\hspace*{16pt}}&

\multicolumn{1}{p{16pt}|}{\hspace*{16pt}}&

\multicolumn{1}{p{16pt}|}{\hspace*{16pt}}\\[-10pt]

R & 2\hphantom{$-$}& 1$-$& 4\hphantom{$-$}& 3$-$\\ 

L & 2$-$& 1\hphantom{$-$}& 4$-$& 3\hphantom{$-$}\\

D & 3\hphantom{$-$}& 4\hphantom{$-$}& 1$-$& 2$-$\\

U & 3$-$& 4$-$& 1\hphantom{$-$}& 2\hphantom{$-$}\\ 

\hline

\end{tabular}

\end{center}

%\vspace*{-6pt}

\end{table}

%\end{floatingfigure}



Ниже приведена последовательность шагов алгоритма поиска соседей для узла:

  \begin{itemize}

\item если {Left}, {Right}, {Top}, {Bottom} равны~0 или 

размеру изображения ($W$, $H$), то в соответствующем направлении соседей 

нет;

\item вычислить соседа для $\mathrm{Node}_n$: в цикле для всех существующих 

направлений (не равных~0) вычислить все значения соседних квадрантов на 

основе $\mathrm{ChNum}$ в соответствии с таблицей;

\item занести значение в~$\mathrm{Env}$.

\end{itemize}



  Конечный автомат возвращает число, обозначающее соседний квадрант в~со\-от\-вет\-ствующем 

направлении. Уровень соседа вычисляется до размера не менее размера 

текущего узла~\cite{14-sor}. Но квадрант может граничить с соседом любого 

размера. В~этом случае необходимо уточнение размера соседей. Поскольку 

$\mathrm{Env}$ содержит идентификатор соседей, необходимо провести направленный 

спуск по дереву и определить существование узла с соответствующим 

идентификатором. Ниже представлен алгоритм поиска существующих соседей.

  

  На вход подается дерево квадрантов $\mathrm{QTree}$ с вычисленным $\mathrm{Env}$. Далее 

выполняются следующие действия:

  \begin{itemize}

\item просмотр значений из $\mathrm{Env}$ и занесение их в стек $\mathrm{strID}$;

\item в цикле осуществить просмотр значений для верхнего значения $\mathrm{strID}$: 

занести его в текущий путь $\mathrm{curWay}$, если его длина равна~0 или узел с 

текущим номером не существует, то закончить дальнейшее рассмотрение 

текущего пути и вернуть существующий путь и указатель на последний 

существующий узел этого пути, иначе занести значение в пройденный путь 

$\mathrm{rtnWay}$, повторить для следующего значения;

\item создать связь из текущего узла в узел с индексом в~$\mathrm{rtnWay}$;

\item сохранить пройденный путь $\mathrm{rtnWay}$ и указатель на узел.

  \end{itemize}

  

  На выходе алгоритма~--- $\mathrm{QTree}$ с вычисленным существующим $\mathrm{Env}$ 

  и~ссылками на соседей. Результатом работы этого алгоритма будет граф, в 

котором имеются переходы от узла к каждому существующему соседу, не 

меньше текущего узла.

  

  Для того чтобы осуществить слияние, необходимо перейти к линейным 

формам квадродерева $\mathrm{QTree}$~--- $\mathrm{LnQTree}$~--- и провести сортировку по 

числу точек в квадранте в порядке возрастания. Это делается для ускорения 

слияния. 

  

  Алгоритм слияния на вход получает дерево в отсортированном линейном 

виде $\mathrm{LnQTree}$ и~предполагает следующую последовательность шагов:

  \begin{itemize}

\item осуществить просмотр массива $\mathrm{LnQTree}$, пока в~нем имеются 

элементы;

\item извлечь крайний элемент массива $\mathrm{LnQTree}$ и~посмотреть массив 

слияния $\mathrm{Segment}$ (содержащий индекс сегмента, которому принадлежит 

квадрант,

число точек и~цвет): если он пуст, то создать первый сегмент, иначе 

создать новый сегмент;



\item пометить узел как посещенный;

\item в цикле осуществить просмотр узлов, смежных с узлами текущего 

сегмента, и~поместить их в~$\mathrm{vNgbrSgmNode}$;

\item вычислить близость цвета сегмента и~узла соседа и,~если она меньше 

порога, добавить этот узел в~сегмент и~отметить узел как посещенный;

\item осуществить просмотр $\mathrm{vNgbrSgmNode}$ и~извлечь узлы, которые были 

посещены;

\item аналогично пройтись по массиву $\mathrm{LnQTree}$ и~удалить посещенные узлы.

  \end{itemize}

  

  На выходе алгоритма~--- массив сегментов, в~котором каждый сегмент 

состоит из узлов, близких по цветовым характеристикам.

  

  Таким образом, каждый узел можно классифицировать и~сопоставить 

  с~определенным сегментом. Принадлежность к~соответствующему сегменту 

определяется исходя из значения, указанного в~$\mathrm{segmentnumber}$, 

используемого в~качестве индекса массива сегментов.

\vspace*{-3pt}



\section{Результаты}



  Результат работы описанного алгоритма приведен на рис.~3. 

\vspace*{-3pt}



\section{Выводы}



  Как видно из представленных результатов, алгоритм сегментации хорошо 

справляется с изображениями, содержащими значительное число мелких 

деталей, потеря которых создает трудности при анализе. При этом фрагменты 

изображения, представляющие собой однородную область, объединяются в 

более крупные сегменты.

  

  Параметры, используемые алгоритмом при сегментации, были установлены 

эмпирически. Возможно, было бы целесообразно выбрать их автоматически 

в~зависимости от типа входного изображения. По мере дальнейшей работы 

возможно улучшение выбора границ и, вероятно, использование 

дополнительных индикаторов, например основанных на текстурах, для 

улучшения результатов сегментации, в частности для сложных изображений. 



\pagebreak



\begin{figure} %fig3

\vspace*{1pt}

\begin{center}

 \mbox{%

 \epsfxsize=128mm 

\epsfbox{man-3.eps}

 }

 \end{center}

\vspace*{-11pt}

\Caption{Результаты работы алгоритма: (\textit{a})~исходное изображение; (\textit{б})~утончение 

контуров; (\textit{в})~уменьшение цветов на изображении; (\textit{г})~разделение; 

(\textit{д})~слияние}

%\vspace*{-2pt}

\end{figure}

     



%\vspace*{-3pt}

  

{\small\frenchspacing

{%\baselineskip=10.8pt

\addcontentsline{toc}{section}{Литература}

\begin{thebibliography}{99}

  

\bibitem{1-sor}

\Au{Stojmenovic M., Montero A.\,S., Nayak~A.} Colour and texture based pyramidal image 

segmentation~/\!\!/  Conference (International) on Audio, Language and Image Processing 

Proceedings.~--- IEEE, 2010. P.~778--786.

\bibitem{2-sor}

\Au{Correa-Tome F.\,E., Sanchez-Yanez R.\,E.} Integral split-and-merge methodology for real-time 

image segmentation~/\!\!/ J.~Electron. Imaging, 2015. Vol.~24. Iss.~1. P.~1--11.

\bibitem{3-sor}

\Au{Волков В.\,Н., Сорокин А.\,И.} Классификация методов сегментации цветных 

изоб\-ра\-же\-ний~/\!\!/ Информационные системы и технологии, 2018. №\,4(108). С.~32--44.

\bibitem{4-sor}

\Au{Fu K.\,S., Mui J.\,K.} A~survey on image segmentation, pattern recognition~/\!\!/ Pattern 

Recogn., 1981. Vol.~13. Iss.~1.  

P.~3--16.

\bibitem{5-sor}

\Au{Ferhat F.\,A., Kerdjidj O., Messaoudi~K., Abdelwahhab~B.} Implementation of Sobel, Prewitt, 

Roberts edge detection on FPGA~/\!\!/ World Congress in Computer Science, Computer Engineering, 

and Applied Computing Proceedings, 2013. P.~1--4.

\bibitem{6-sor}

\Au{Lucchese L., Mitray S.\,K.} Color image segmentation: A~state-of-the-art survey~/\!\!/ 

P.~Indian National Science Academy, 2001. Vol.~67. Iss.~2. P.~207--221.

\bibitem{7-sor}

\Au{Verges-Llahi J.} Color constancy and image segmentation techniques for applications to mobile 

robotics.~--- Barcelona, Spain: Polytechnic University of Catalonia, 2005.  PhD Thesis. 23~p.

\bibitem{8-sor}

\Au{Lorenzo-Navarro J., Hern$\acute{\mbox{a}}$ndez-Tejera M.} Image segmentation using a 

modified split and merge technique~/\!\!/ Cybernet. Syst., 2007. Vol.~25. Iss.~1. P.~137--162.

\bibitem{9-sor}

\Au{Aizawa K., Motomura K., Kimura~S., Kadowaki~R., Fan~J.} Constant time neighbor finding in 

quadtrees: An experimental result~/\!\!/ 3rd  Symposium (International) on Communications, Control and 

Signal Processing Proceedings.~--- IEEE, 2008. P.~505--510.



\bibitem{10-sor}

\Au{Kelkar D., Gupta S.} Improved quadtree method for split merge image segmentation~/\!\!/ 1st 

Conference (International) on Emerging Trends in Engineering and Technology Proceedings.~---

IEEE, 2008. P.~44--47.

\bibitem{11-sor}

\Au{Сорокин А.\,И., Волков В.\,Н. Шульдешова~О.\,В.} Синтетический алгоритм обнаружения 

контуров объектов на основе  

RGB-ко\-ор\-ди\-нат~/\!\!/ Информационные системы и технологии, 2019. №\,2(112). С.~27--

34.

\bibitem{12-sor}

\Au{Kwon J.\,S., Woong J., Kang E.\,K.} An enhanced thinning algorithm using parallel processing~/\!\!/ 

Conference (International) on Image Processing Proceedings.~--- IEEE, 2001. Vol.~3. P.~752--755.

\bibitem{13-sor}

\Au{Lagani$\grave{\mbox{e}}$re R.} OpenCV~2 computer vision application programming 

cookbook.~--- Packt Publishing, 2011. 306~p.

\bibitem{14-sor}

\Au{Yoder R., Bloniarz P.} A~practical algorithm for computing neighbors in quadtrees, octrees, and 

hyperoctrees~/\!\!/  Conference (International) on Modeling, Simulation \&~Visualization Methods

Proceedings.~--- 

CSREA Press, 2006. P.~249--255.

\end{thebibliography}

} }





%\vspace*{-3pt}



\hfill{\small\textit{Поступила в~редакцию 14.09.20}}









\vspace*{8pt}



\hrule



\vspace*{2pt}



%\newpage



%\vspace*{-28pt}



\hrule



\vspace*{8pt}



\def\tit{QUADTREE BASED COLOR IMAGE SEGMENTATION METHOD}



\def\titkol{Quadtree based color image segmentation method}



\def\aut{Yu.\,A.~Maniakov and A.\,I.~Sorokin}



\def\autkol{Yu.\,A.~Maniakov and A.\,I.~Sorokin}



\titel{\tit}{\aut}{\autkol}{\titkol}



\vspace*{-5pt}





\noindent

Orel Branch of the Federal Research Center ``Computer Science and Control''

of the Russian Academy of Sciences, 

137~Moskovskoe Shosse, Orel 302025, Russian Federation





\def\leftfootline{\small{\textbf{\thepage}

\hfill Sistemy i Sredstva Informatiki~---

Systems and Means of Informatics 2020 vol~30 no\,4}

}%

 \def\rightfootline{\small{Sistemy i Sredstva Informatiki~---

 Systems and Means of Informatics 2020 vol~30 no\,4

\hfill \textbf{\thepage}}}

  

  

\vspace*{6pt}





\Abste{ The paper presents a color image segmentation method and an algorithm based on quadtree. The 

proposed method consists of several steps. First of them is border detection based on three-channel color and 

two masks. Then, the authors apply the thinning algorithm to decrease the area of the found boundary. The 

segmentation algorithm is divided into two parts. In the first part, the image is divided into segments as much as 

possible. In the second part, the segments union algorithm uses the finding neighbor's ID

based on the FSM table and applies a~link to ID to create a graph. The results of color images segmentation obtained on the basis of the 

described algorithm are presented.}



%%\pagebreak



\KWE{color image; segmentation; quadtree; edge; border; thinning; color reduction; split; merge; pixel}



\DOI{10.14357/08696527200410}





%\vspace*{-24pt}



%\Ack

%\noindent

 



\vspace*{-12pt}



\renewcommand{\bibname}{\protect\rmfamily References}

%\renewcommand{\bibname}{\large\protect\rm References}



{\small\frenchspacing

{%\baselineskip=10.8pt

\addcontentsline{toc}{section}{References}

\begin{thebibliography}{99}



\bibitem{1-sor-1}

\Aue{Stojmenovic, M., A. S. Montero, and A.~Nayak.} 2010. Colour and texture based pyramidal image 

segmentation. \textit{Conference (International) on Audio, Language and Image Processing Proceedings}. 

IEEE. 778--786.

\bibitem{2-sor-1}

\Aue{Correa-Tome, F.\,E., and R.\,E. Sanchez-Yanez.} 2015. Integral split-and-merge methodology for 

real-time image segmentation. \textit{J.~Electron. Imaging} 24(1):1--11.

\bibitem{3-sor-1}

\Aue{Volkov, V.\,N., and A.\,I.~Sorokin.} 2018. Klassifikatsiya metodov segmentatsii tsvetnykh izobrazheniy 

[Classification of methods of segmentation color image]. \textit{Informatsionnye sistemy i~tekhnologii} 

[Information Systems and Technologies] 4(108):32--44.

\bibitem{4-sor-1}

\Aue{Fu, K.\,S., and J.\,K. Mui.} 1981. A~survey on image segmentation, pattern recognition. \textit{Pattern 

Recogn.} 13(1):3--16.

\bibitem{5-sor-1}

\Aue{Ferhat, F.\,A., O. Kerdjidj, K. Messaoudi, and B.~Abdelwahhab.} 2013. Implementation of Sobel, Prewitt, 

Roberts edge detection on FPGA. \textit{World Congress in Computer Science, Computer Engineering, and 

Applied Computing Proceedings}. 1--4.

\bibitem{6-sor-1}

\Aue{Lucchese, L., and S.\,K. Mitray.} 2001. Color image segmentation: A~state-of-the-art survey. 

\textit{P.~Indian National Science Academy} 67(2):207--221.

\bibitem{7-sor-1}

\Aue{Verges-Llahi, J.} 2005. Color constancy and image segmentation techniques for applications to mobile 

robotics. Barcelona, Spain: Polytechnic University of Catalonia. PhD Thesis.  23~p.

\bibitem{8-sor-1}

\Aue{Lorenzo-Navarro, J., and M. Hern$\acute{\mbox{a}}$ndez-Tejera.} 2007. Image segmentation using 

a~modified split and merge technique. \textit{Cybernet. Syst.} 25(1):137--162.

\bibitem{9-sor-1}

\Aue{Aizawa, K., K. Motomura, S. Kimura, R.~Kadowaki, and J.~Fan.} 2008. Constant time neighbor finding in 

quadtrees: An experimental result. \textit{3rd Symposium (International) on Communications, Control and 

Signal Processing Proceedings}. IEEE. 505--510.

\bibitem{10-sor-1}

\Aue{Kelkar, D., and S. Gupta.} 2008. Improved quadtree method for split merge image segmentation. 

\textit{1st Conference (International) on Emerging Trends in Engineering and Technology Proceedings}. 

IEEE. 44--47.

\bibitem{11-sor-1}

\Aue{Sorokin, A.\,I., V.\,N. Volkov, and O.\,V.~Shuldeshova.} 2019. Sinteticheskiy algoritm obnaruzheniya 

konturov ob''ektov na osnove RGB-koordinat [Synthetic algorithm for detection of object contours on the raster 

images based on RGB coordinates]. \textit{In\-for\-ma\-tsi\-on\-nye sistemy i~tekhnologii} [Information Systems and 

Technologies] 2(112):27--34.

\bibitem{12-sor-1}

\Aue{Kwon, J.\,S., J. Woong, and E.\,K.~Kang.} 2001. An enhanced thinning algorithm using parallel 

processing. \textit{Conference (International) on Image Processing Proceedings}. IEEE. 752--755.

\bibitem{13-sor-1}

\Aue{Lagani$\grave{\mbox{e}}$re, R.} 2011. \textit{OpenCV 2 computer vision application programming cookbook}. Packt 

Publishing. 306~p.

\bibitem{14-sor-1}

\Aue{Yoder, R., and P. Bloniarz.} 2006. A~practical algorithm for computing neighbors in quadtrees, octrees, 

and hyperoctrees. \textit{Conference (International) on Modeling, Simulation \& Visualization Methods 

Proceedings.} CSREA Press. 249--255.

\end{thebibliography}

} }



%\vspace*{-3pt}



\hfill{\small\textit{Received September 14, 2020}} 



%\vspace*{-8pt}



\Contr



\noindent

\textbf{Maniakov Yuri A.} (b.\ 1984)~--- senior scientist, Orel Branch of the Federal Research Center 

``Computer Science and Control'' of the Russian Academy of Sciences, 137~Moskovskoe Shosse, Orel 302025, 

Russian Federation; \mbox{maniakov\_yuri@mail.ru}



\vspace*{5pt}







\noindent

\textbf{Sorokin Andrey I.} (b.\ 1987)~--- junior scientist, Orel Branch of the Federal Research Center 

``Computer Science and Control'' of the Russian Academy of Sciences, 137 Moskovskoe Shosse, Orel 302025, 

Russian Federation; \mbox{webdi@mail.ru}



\label{end\stat}



\renewcommand{\bibname}{\protect\rm Литература} 

